% Anexo: herramientas de línea de comandos (ajuste headless).
% Se incluye con \input desde main.tex bajo un \chapter, así que empieza en \section.

Every fit that the GUI performs can also be run \emph{headless} from the
command line, with the \textbf{same numerical engine}
(\code{core.fit\_engine} and \code{mossbauer\_distribution}): results are
identical to those obtained interactively. Two scripts cover the two fitting
worlds:

\begin{itemize}[nosep]
  \item \file{mossbauer\_fit\_cli.py} --- \textbf{discrete} fits
        (singlet / doublet / sextet / relaxation models) driven by a session
        template, including Monte-Carlo bootstrap errors, profile-likelihood
        intervals and warm-start batch series.
  \item \file{fit\_bhf\_distribution\_cli.py} --- \textbf{hyperfine
        distribution} fits: $P(\BHF)$ or $P(\dEQ)$, every shape
        (histogram, Gaussian, VBF, binomial), every regularizer
        (Tikhonov, TV, MaxEnt), Voigt kernels, field correlations and the 2D
        distribution $P(\BHF,\dEQ)$.
\end{itemize}

Neither script needs a display: they run on a cluster, inside a batch queue,
or from another program. Both return exit code $0$ on success, which makes
them easy to chain in shell scripts.

\section{Discrete fits: \file{mossbauer\_fit\_cli.py}}
\label{sec:cli-discrete}

The workflow mirrors the GUI: prepare the model interactively once, save it as
a session (\menu{File $\triangleright$ Save session\ldots}), and use that JSON
file as a \textbf{template} for any number of spectra:

\begin{lstlisting}[language=bash]
python mossbauer_fit_cli.py \
    --template template_alphaFe.json \
    --spectrum sample.adt \
    --out sample.fit.json
\end{lstlisting}

The template carries the full model definition: component kinds and starting
values, fixed/free flags, constraints, line profile (Lorentzian/Voigt),
likelihood (Gaussian/Poisson), robust loss, absorber model, drive waveform
(triangular/sine), multistart count and the calibration flags. The folding
centre is \emph{not} taken from the template --- it is detected for each
spectrum, exactly as the GUI does. The output JSON has the same structure as
a GUI session (it can be loaded back with
\menu{File $\triangleright$ Load session\ldots}) plus a
\code{batch\_fit\_result} block with the fitted values, errors and statistics
($\chi^2$, reduced $\chi^2$, AIC, BIC).

\subsection{Options}

\begin{center}
\begin{tabular}{@{}llp{7.6cm}@{}}
\toprule
Option & Argument & Meaning \\
\midrule
\code{-{}-template} & \file{file.json} & Session template (required). \\
\code{-{}-spectrum} & one or more files & Spectrum(s) \file{.ws5}/\file{.adt}
  to fit. With several files the series runs with sequential warm-start
  (below). \\
\code{-{}-out} & \file{file.json} & Output file (required). \\
\code{-{}-vmax} & mm/s & Overrides the template $v_{\max}$ (useful when the
  calibration changes between series). \\
\code{-{}-bootstrap} & $N$ & Monte-Carlo bootstrap errors with $N$ replicas
  (single spectrum only). Adds a \code{bootstrap} section with the per-parameter
  $\sigma_{\mathrm{MC}}$. \\
\code{-{}-profile-likelihood} & --- & Asymmetric $1\sigma/2\sigma$ intervals by
  profile likelihood ($\Delta\chi^2=1$ and $\Delta\chi^2=4$ crossings; see the
  mathematical appendix). Adds a \code{profile\_likelihood} section. \\
\code{-{}-quiet} & --- & Suppress the summary printed to stdout. \\
\bottomrule
\end{tabular}
\end{center}

\subsection{Advanced errors}

With \code{-{}-bootstrap} $N$ the base fit is repeated over $N$ synthetic
spectra (real Poisson counts when the normalization factor is available) and
the sample standard deviation of each free parameter is reported next to the
covariance error:

\begin{lstlisting}[language=bash]
python mossbauer_fit_cli.py --template T.json --spectrum a.adt \
    --out a.fit.json --bootstrap 100
#  s1_bhf = 33.0446 +- 0.005988  [sigma_MC 0.006113]
\end{lstlisting}

With \code{-{}-profile-likelihood} each free parameter is scanned around its
optimum while re-fitting the rest, and the $\Delta\chi^2=1$ / $\Delta\chi^2=4$
crossings give asymmetric intervals that remain valid when the parabolic
(covariance) approximation fails:

\begin{lstlisting}[language=bash]
#  s1_bhf = 33.0446 +- 0.005988  [1-sigma profile -0.00607/+0.006066]
\end{lstlisting}

\subsection{Warm-start series}

When several spectra are given (a temperature or composition series), each
fit starts from the converged values of the previous one --- the same
sequential warm-start as the GUI batch tool. The output is a single JSON with
one row per file (\code{status}, \code{values}, \code{errors}, \code{stats});
individual failures do not stop the series:

\begin{lstlisting}[language=bash]
python mossbauer_fit_cli.py --template T.json \
    --spectrum serie_010K.adt serie_050K.adt serie_100K.adt \
    --out serie.batch.json
\end{lstlisting}

For fully independent fits (no warm-start) simply call the script once per
file from the shell.

\section{Distribution fits: \file{fit\_bhf\_distribution\_cli.py}}
\label{sec:cli-dist}

The distribution CLI reads a raw \file{.ws5}/\file{.adt} file, folds it
(NORMOS-style, reusing the \file{.RES}/\file{.JOB} sidecars when present) and
fits a hyperfine distribution with the same engine as the GUI panel of
chapter~\ref{ch:distribuciones}:

\begin{lstlisting}[language=bash]
# Classic P(BHF) histogram (Tikhonov):
python fit_bhf_distribution_cli.py sample.adt --alpha 0.01 --nbins 50 --plot

# P(DEQ) doublet distribution:
python fit_bhf_distribution_cli.py sample.adt --variable quad --bmax 3

# VBF with two Gaussian components; MaxEnt histogram; 2D distribution:
python fit_bhf_distribution_cli.py sample.adt --shape vbf --vbf-components 2
python fit_bhf_distribution_cli.py sample.adt --reg-mode maxent
python fit_bhf_distribution_cli.py sample.adt --dist-2d --qmin -1 --qmax 1
\end{lstlisting}

\subsection{Model options}

\begin{center}
\begin{tabular}{@{}llp{7.2cm}@{}}
\toprule
Option & Values & Meaning \\
\midrule
\code{-{}-variable} & \code{bhf} $|$ \code{quad} & Distributed variable:
  $P(\BHF)$ (default) or $P(\dEQ)$ with the field fixed at
  \code{-{}-fixed-bhf} ($0$~T $\Rightarrow$ doublet kernel). \\
\code{-{}-shape} & \code{histograma} $|$ \code{gaussiana} $|$ \code{vbf} $|$
  \code{binomial} & Distribution shape (chapter~\ref{ch:distribuciones}). \code{gaussiana}
  is the VBF with $N=1$ and a Lorentzian line. \\
\code{-{}-vbf-components} & $N$ & Number of Gaussians of the VBF. \\
\code{-{}-reg-mode} & \code{tikhonov} $|$ \code{tv} $|$ \code{maxent} &
  Histogram regularizer. \\
\code{-{}-profile} & \code{Lorentziana} $|$ \code{Voigt} & Kernel line profile
  (\code{-{}-voigt-sigma} sets the instrumental Gaussian width, mm/s). \\
\code{-{}-delta-slope}, \code{-{}-quad-slope} & slope & Linear correlations
  $\delta(H)$ and $\dEQ(H)$ of appendix~\ref{ap:mates}; $0$ = classic kernel. \\
\code{-{}-bmin}, \code{-{}-bmax}, \code{-{}-nbins} & grid & Grid of the
  distributed variable (T for \code{bhf}, mm/s for \code{quad}). \\
\code{-{}-delta}, \code{-{}-quad}, \code{-{}-gamma} & values & Fixed kernel
  parameters; defaults are read from the NORMOS sidecars when available. \\
\code{-{}-sharp-bhf} & T (repeatable) & Adds crystalline (sharp) sextets on
  top of the distribution, e.g.\ \code{-{}-sharp-bhf 33 -{}-sharp-bhf 49}. \\
\code{-{}-dist-2d} & --- & Two-dimensional $P(\BHF,\dEQ)$
  (appendix~\ref{ap:dist-2d}), with \code{-{}-qmin/-{}-qmax/-{}-nbins-quad}
  and \code{-{}-alpha-quad}. \\
\code{-{}-scan-alpha} & --- & Scans $\alpha$ (histogram only) and writes a
  simple L-curve with a suggested corner value. \\
\bottomrule
\end{tabular}
\end{center}

\subsection{Output files}

For 1D fits three files are written next to the input (or under
\code{-{}-out-prefix}): \file{*\_spectrum.dat} (velocity, data, fit,
residual, folded counts), \file{*\_distribution.dat} (grid, $P$, normalized
$P$) and \file{*\_summary.json} (all scalar parameters and metrics; for the
VBF it includes the fitted $(A_j,\mu_j,\sigma_j)$ triplets). With
\code{-{}-plot} a PNG with the spectrum and the distribution is added.

In 2D mode the distribution is written as a matrix
(\file{*\_distribution2d.dat}) with its two axes in
\file{*\_bhf\_axis.dat} / \file{*\_quad\_axis.dat}, and the summary carries
the moments $\langle\BHF\rangle$, $\langle\dEQ\rangle$, widths and the
$\BHF$--$\dEQ$ correlation.
